\chapter{Related Work, Method Families, and the Gap ORIUS Addresses}
\label{ch:related-work-universal}
\section{{Conformal and distribution-free uncertainty}}
Distribution-free predictive inference provides one of the few practical ways to place
finite-sample uncertainty guards around learned models without assuming a perfectly
specified parametric error law \cite{vovk2005algorithmic,lei2018distribution,romano2019conformalized,gibbs2021adaptive,angelopoulos2023gentle,barber2023beyond,zaffran2022adaptive,stankeviciute2021conformal,xu2021dynamic,sesia2021comparison,bates2021distribution}.  ORIUS adopts this family not because
conformal prediction alone solves runtime safety, but because it supplies a principled
outer shell around the unobserved state.  What this family solves is calibration of
predictive uncertainty.  What it does \emph{not} solve is the runtime legality of the
action that will be executed after telemetry degrades.  An interval can remain
statistically respectable while the controller still reasons over the wrong physical
state.  The gap ORIUS addresses here is the missing runtime bridge between nominal
predictive coverage and actuation legality under degraded telemetry.

\section{{Runtime assurance, supervisory safety, and shields}}
Runtime assurance and shielding treat safety as a supervisory execution discipline rather
than a property delegated to the nominal controller \cite{sha2001using,sha1998simplex,desai2019soter,bloem2015shield,bartocci2018introduction,havelund2004runtime,leucker2009brief,schierman2015run,bak2011runtime}.  Simplex-style
architectures, runtime monitors, and shield synthesis all contribute the idea that
proposed actions may be intercepted, replaced, or downgraded at runtime.  The gap ORIUS
addresses here is that most runtime assurance architectures begin \emph{after} the
nominal controller has already proposed an action.  They often assume the relevant state
estimate is coherent enough for the safety supervisor to interpret.  ORIUS instead treats
\emph{observation degradation itself} as the mechanism that makes action legality
uncertain in the first place, which means the assurance layer must reason jointly about
reliability, uncertainty inflation, repair geometry, and certificate semantics.

\section{{Safety filters, barrier methods, and robust control}}
Control barrier functions, robust MPC, reachable sets, predictive safety filters, and
constrained control all provide machinery for turning uncertainty into admissible action
sets \cite{ames2019control,ames2014cbf,wabersich2021predictive,nguyen2016exponential,wang2017safety,fisac2019general,mayne2005robust,langson2004tubes,rawlings2017mpc,ben2009robust,camacho2013model}.  Their shared limitation, from the perspective of this manuscript, is
not lack of mathematical rigor.  It is that they are usually written at the level of one
plant, one controller class, or one constraint geometry.  The missing object is a typed
interface that binds telemetry reliability, uncertainty inflation, repair, fallback, and
certificate emission into a reusable cross-domain contract.  ORIUS is written to supply
that interface rather than to displace the underlying control tools.

\section{{Drift, anomaly detection, and cyber-physical degradation}}
Change detection, anomaly detection, CPS security, and fault diagnosis provide the
language needed to recognize when the observation channel is no longer trustworthy
\cite{rajkumar2010cps,derler2012modeling,baheti2011cyberphysical,pasqualetti2013attack,teixeira2015secure,basseville1993detection,page1954continuous,page1955test,hinkley1971inference,liu2008isolation,breunig2000lof}.  ORIUS depends on that language, but departs from it by treating degraded
observation not only as a detection problem but as an \emph{action semantics} problem.
Most detection papers terminate at an alarm, a score, or a classifier.  ORIUS asks the
next question that matters operationally: once trust in the observation is degraded, how
must the admissible action set, fallback behavior, and certificate payload change before
the system is allowed to act again?

\section{{Applied domain literatures}}
The applied rows of the book draw from multiple domain literatures rather than from a
single benchmark tradition.  Energy and smart-grid control literature motivates the
reference witness; autonomous-vehicle and robotics literature motivates the need for
runtime repair under perception uncertainty; industrial process-control literature
provides real operational envelope constraints; healthcare monitoring shows how degraded
observation can suppress necessary intervention; and navigation and aerospace show where
the present closure boundary still ends.  These families collectively justify a universal
framing, but they also make the insufficiency of prior work visible: every family solves
an adjacent part of the problem, and none alone supplies a universal runtime grammar for
degraded observation, repaired actuation, and auditable safety certificates.  That is the
gap ORIUS is designed to close.
